\chapter*{\textit{Abstract}}
\addcontentsline{toc}{chapter}{\textit{Abstract}}

{
    \setstretch{1}
    
    \begin{center}
        % --- Penulis dan Email Korespondensi ---
        \namaMahasiswa \\
        \texttt{\emailref{\emailKorespondensi}}
        
        % --- Judul ---
        \textit{\textbf{\MakeUppercase{\judulENG}}}
        
        % --- Informasi Institusi & Tahun ---
        \textit{\utDaerah, Faculty of \fakultasENG, \programStudiENG\ Study Program, \tahun}
    \end{center}
    
    \vspace{0.4cm}
    
    \begin{otherlanguage*}{american}
    \noindent
    {\itshape
    %------------------------ BATAS AMAN ------------------------
    % --- Isi Teks Abstrak Bahasa Inggris ---
    %
    %     Tulis isi teks abstrak Anda dengan maksimal 250 kata
    %
    The abstract is written as text that presents the document's essence independently without referencing tables, figures, or bibliographies. The composition begins with the background or research rationale that underlies the importance of the studied or developed topic. Once the urgency of the problem is briefly stated, the next section must clearly formulate the primary objective to be achieved in the study. The next crucial component is the explanation of the method or problem-solving approach used by the researcher. This methodological description must encompass the research design, location, subjects or data sources, data collection techniques, instruments used, and data analysis techniques applied to process information. After the methods are outlined, the structure shifts to presenting the main findings obtained based on the objective analysis of the research object. The final part of the abstract body is closed with conclusions that directly answer the initial objective, along with relevant operational suggestions for future scientific development or practical policies. The entire description must be arranged concisely, clearly, and systematically using single spacing, and the total length must not exceed the maximum limit of two hundred and fifty words. Choose representative vocabulary to be used as keywords at the bottom of the abstract, then arrange these words alphabetically from A to Z with a maximum of six keywords separated by commas. Ensure that the abstract heading format complies with the faculty handbook before the document is ready to be compiled into the final printed version of your thesis.
    %------------------------ BATAS AMAN ------------------------
    }
    
    \vspace{0.6cm}
    
    \noindent
    %------------------------ BATAS AMAN ------------------------
    % --- Kata Kunci ---
    %
    %     Maksimal 6 kata, diurutkan dari A sampai Z
    %
    \textbf{Keywords:} Analysis, E-Commerce, Kelayakan, Kualitatif, Skripsi, UMKM.
    %------------------------ BATAS AMAN ------------------------
    \end{otherlanguage*}
}

\newpage